% Source: https://artofproblemsolving.com/community/c3h1850899_easy_recursion_function_problem
% Stable identity: aops-topic-1850899-q01
\textbf{Small values.} Using $f(x)=x-2$ for $x>5$:
\[
\begin{aligned}
f(5)&=5-f(f(8))=5-f(6)=5-4=1,\\
f(4)&=5-f(f(7))=5-f(5)=5-1=4,\\
f(3)&=5-f(f(6))=5-f(4)=5-4=1,\\
f(1)&=5-f(f(4))=5-f(4)=5-4=1,\\
f(2)&=5-f(f(5))=5-f(1)=5-1=4.
\end{aligned}
\]
So \textbf{(a)} $f(1)=\boxed{1}$. Note that $f(1)=1$ and $f(4)=4$.

\textbf{(c)} We show by downward induction that for every integer $x\le5$, $f(x)$ is determined by the rules and
\[
f(x)=\begin{cases}1,&x\text{ odd},\\4,&x\text{ even}.\end{cases}
\]
This holds for $x=1,\dots,5$ by the computation above. Let $x\le0$ and suppose it holds for all integers $y$ with $x<y\le5$. Then $x+3\le3$, so $f(x+3)$ is already known and equals $1$ or $4$; since $f(1)=1$ and $f(4)=4$, we get $f(f(x+3))=f(x+3)$ and
\[
f(x)=5-f(x+3).
\]
The number $x+3$ has the opposite parity to $x$. If $x$ is odd, $f(x+3)=4$ and $f(x)=1$; if $x$ is even, $f(x+3)=1$ and $f(x)=4$. This completes the induction. Every value was forced, so $f$ is uniquely determined, and conversely this $f$ satisfies the rules (the same computations check them). The closed form is
\[
f(x)=\begin{cases}x-2,&x>5,\\1,&x\le5,\ x\text{ odd},\\4,&x\le5,\ x\text{ even}.\end{cases}
\]

\textbf{(b)} Since $-10$ is even and at most $5$, $f(-10)=\boxed{4}$.
